\begin{resumo}[Abstract]
 \begin{otherlanguage*}{english}
   This monograph presents the application of complex network analysis to historical Olympic data, modeling sports competitions through directed and weighted graphs that reveal latent competitive structures. The work analyzes 9,350 medalist athletes across six sports (Swimming, Basketball, Football, Athletics, Judo, Boxing) spanning 125 years of Olympic history (1896-2021), constructing networks where athletes constitute nodes and competitive relationships (confrontations in shared podiums) form edges weighted by medal hierarchy. Through PageRank adapted to sports contexts, centrality metrics, and community detection via Louvain algorithm, the study identifies performance hierarchies, structural clusters, and temporal patterns not evident in traditional medal-count analyses. Results reveal systematic differences between modalities: individual sports produce highly modular networks with segregated communities, while team sports generate densely interconnected structures with low modularity. Multidimensional community characterization through temporal entropy, geographic concentration, and dominance indices enables deep understanding of Olympic competitive dynamics. An interactive Streamlit dashboard facilitates results exploration and findings validation.

   \vspace{\onelineskip}

   \noindent
   \textbf{Keywords}: complex networks, sports performance analysis, PageRank, community detection, Olympic sport, graph modeling, centrality metrics, Louvain algorithm.
 \end{otherlanguage*}
\end{resumo}